% Figure: Maekawa barrier insertion loss versus Fresnel number N, the
% empirical design curve for a thin rigid barrier in the shadow zone.
\begin{figure}[htbp]
\centering
\begin{tikzpicture}
\begin{axis}[
  width=12cm, height=7.5cm,
  xlabel={Fresnel number $N = 2\delta/\lambda$ (log scale)},
  ylabel={barrier insertion loss (dB)},
  xmode=log, log basis x=10,
  xmin=0.05, xmax=30, ymin=0, ymax=26,
  axis lines=left,
  domain=0.05:30, samples=120,
  legend pos=south east,
  legend cell align=left,
  legend style={font=\scriptsize},
  tick label style={font=\scriptsize},
  label style={font=\small},
  grid=major, grid style={enggray!20},
]
% Maekawa: IL = 10 log10(3 + 20 N) for N > 0, clamped at 5 dB minimum near N=0
\addplot[engblue, very thick] {10*log10(3 + 20*x)};
\addlegendentry{$\text{IL}=10\log_{10}(3+20N)$}
% the 5 dB practical floor (grazing incidence)
\addplot[enggray, thin, dashed] {5};
\node[enggray, font=\scriptsize, anchor=south west] at (axis cs:0.055,5.2) {$\approx 5$\,dB grazing floor};
% mark the worked design point N = 1.6 -> IL = 10 log10(3+32)=15.4 dB
\addplot[engblack, only marks, mark=*, mark size=2pt] coordinates {(1.6,15.4)};
\node[engblack, font=\scriptsize, anchor=north west] at (axis cs:1.7,15.0) {worked point $N=1.6$};
\end{axis}
\end{tikzpicture}
\caption[Maekawa barrier insertion-loss curve]{The Maekawa empirical
insertion loss of a thin rigid barrier, $\text{IL}\approx
10\log_{10}(3+20N)$, against the Fresnel number $N=2\delta/\lambda$, where
$\delta$ is the path-length detour over the barrier top and $\lambda$ the
wavelength. The loss climbs about $3\,\si{\decibel}$ per doubling of $N$,
so a barrier that is tall or close to the source or the receiver, or that
blocks a high-frequency source, performs best. Because $N$ carries
$1/\lambda$, the same barrier delivers far more attenuation to a
$2\,\si{\kilo\hertz}$ source than to a $125\,\si{\hertz}$ one, which is
why roadside barriers tame tyre and engine whine yet barely touch the
low-frequency rumble of heavy trucks. The marked point is the worked
highway-barrier example, $N=1.6$, giving about $15\,\si{\decibel}$.}
\label{fig:vol03ch07:barrier-maekawa}
\end{figure}
